\documentclass[11pt, a4paper]{article}

% ── fonts & encoding ────────────────────────────────────────────────────────
\usepackage[T1]{fontenc}
\usepackage[utf8]{inputenc}
\usepackage{mathpazo}          % Palatino for text + math
\usepackage{eulervm}           % Euler math font (AMS-style)
\linespread{1.08}

% ── math ────────────────────────────────────────────────────────────────────
\usepackage{amsmath, amssymb, amsthm}
\usepackage{mathtools}

% ── layout ──────────────────────────────────────────────────────────────────
\usepackage[margin=2.5cm]{geometry}
\usepackage{microtype}
\usepackage{parskip}

% ── figures & tables ────────────────────────────────────────────────────────
\usepackage{graphicx}
\usepackage{float}
\usepackage{caption}
\usepackage{subcaption}
\usepackage{booktabs}

% ── hyperlinks ───────────────────────────────────────────────────────────────
\usepackage[colorlinks=true, linkcolor=black, citecolor=black, urlcolor=blue]{hyperref}

% ── theorem environments ─────────────────────────────────────────────────────
\theoremstyle{plain}
\newtheorem{theorem}{Theorem}[section]
\newtheorem{proposition}[theorem]{Proposition}
\newtheorem{corollary}[theorem]{Corollary}
\theoremstyle{definition}
\newtheorem{definition}[theorem]{Definition}
\newtheorem{remark}[theorem]{Remark}

% ── notation ─────────────────────────────────────────────────────────────────
\newcommand{\E}{\mathbb{E}}
\newcommand{\R}{\mathbb{R}}
\newcommand{\C}{\mathcal{C}}
\newcommand{\pbar}{\bar{p}}
\newcommand{\xstar}{x^*}

% ═══════════════════════════════════════════════════════════════════════════
\begin{document}
% ═══════════════════════════════════════════════════════════════════════════

\title{%
  \large\textsc{Stochastic Battery Valuation}\\[0.4em]
  \LARGE\textbf{The Jensen Gap, Polytope Vertices,\\[0.2em]
  and the Convexity of the Daily Value Function}\\[0.5em]
  \normalsize Technical Note
}
\author{}
\date{\today}
\maketitle
\tableofcontents
\newpage

% ────────────────────────────────────────────────────────────────────────────
\section{Setting}
% ────────────────────────────────────────────────────────────────────────────

We consider a grid-scale electrical storage asset optimised over a
discrete set of trading days $d = 0, 1, \ldots, D{-}1$.
On each day $d$ the asset observes 24 hourly day-ahead prices
$p \in \R^{24}$ and chooses a dispatch schedule
$x = (e^{\mathrm{in}}, e^{\mathrm{out}}) \in \R^{48}$
subject to the physical constraints of the storage device.

The \emph{feasible set} for a given start-of-day energy level
$s_0 \in [s_{\min}, s_{\max}]$ and end-of-day level
$s_1 \in [s_{\min}, s_{\max}]$ is the convex polytope
\begin{equation}
  \C(s_0, s_1)
  \;=\;
  \Bigl\{
    (e^{\mathrm{in}}, e^{\mathrm{out}}) \in \R^{48}
    \;\Big|\;
    \text{capacity, energy-path, and balance constraints}
  \Bigr\},
  \label{eq:feasible}
\end{equation}
where the balance constraint enforces
$\sum_h e^{\mathrm{in}}_h - \sum_h e^{\mathrm{out}}_h = s_1 - s_0$.
A precise statement is given in Appendix~\ref{app:lp}.

The \emph{one-day value function} is
\begin{equation}
  F(p; s_0, s_1)
  \;=\;
  \max_{x \in \C(s_0, s_1)}
  \langle c(p),\, x \rangle,
  \qquad
  c(p)_h
  =
  \begin{pmatrix}
    -p_h / \eta_c \\
    p_h \cdot \eta_d
  \end{pmatrix},
  \label{eq:F}
\end{equation}
where $\eta_c$ and $\eta_d$ are the charging and discharging efficiencies.
Since we discovered empirically that the optimal end-of-day energy level
tends to equal $s_{\min}$ (the battery is fully dispatched within each day),
we simplify below by fixing $s_0 = s_1 = s_{\min}$ and write
$F(p) \equiv F(p; s_{\min}, s_{\min})$.

% ────────────────────────────────────────────────────────────────────────────
\section{Convexity of \texorpdfstring{$F$}{F}}
% ────────────────────────────────────────────────────────────────────────────

\begin{proposition}
$F : \R^{24} \to \R$ is \emph{convex} and \emph{piecewise linear}.
\end{proposition}
\begin{proof}
$F(p) = \sup_{x \in \C} \langle c(p), x \rangle$ is the pointwise
supremum of the family of linear functions $\{p \mapsto \langle c(p), x \rangle\}_{x \in \C}$.
Since $\C$ is a polytope (finitely many vertices), the supremum is over
a finite family of linear functions and is therefore convex and piecewise linear.
\end{proof}

\begin{remark}
The kink points of $F$ — where the optimal vertex changes — correspond
exactly to the \emph{vertex switches} visible in the cross-section plots
(orange dashed lines in Figure~\ref{fig:cross}).
Between consecutive kinks, $F$ is perfectly linear and the gradient equals
the optimal dispatch $\xstar(p)$ evaluated at any interior point.
\end{remark}

% ────────────────────────────────────────────────────────────────────────────
\section{The Jensen Gap}
% ────────────────────────────────────────────────────────────────────────────

\subsection{Definition}

Let $p_1, \ldots, p_N \in \R^{24}$ be $N$ Monte Carlo price scenarios for
day~$d$.  Define
\begin{equation}
  V^{\mathrm{full}}
  \;=\;
  \frac{1}{N} \sum_{i=1}^N F(p_i)
  \;=\;
  \E\bigl[F(p_i)\bigr],
  \qquad
  V^{\mathrm{intr}}
  \;=\;
  F\!\left(\frac{1}{N}\sum_{i=1}^N p_i\right)
  \;=\;
  F(\pbar).
\end{equation}

\begin{definition}[Jensen gap]
The \emph{Jensen gap} for day~$d$ is
\begin{equation}
  \Delta
  \;=\;
  V^{\mathrm{full}} - V^{\mathrm{intr}}
  \;=\;
  \E[F(p_i)] - F(\pbar)
  \;\geq\; 0.
  \label{eq:gap}
\end{equation}
It equals the \emph{extrinsic value} for that day.
\end{definition}

The non-negativity follows immediately from Jensen's inequality applied
to the convex function $F$.

\subsection{Decomposition via individual contributions}

Define the individual contribution of scenario $i$ as
\begin{equation}
  \delta_i \;=\; F(p_i) - F(\pbar).
\end{equation}
Then $\Delta = \frac{1}{N}\sum_i \delta_i = \E[\delta_i]$.
Figure~\ref{fig:hist} shows the empirical distribution of $\delta_i$:
its mean is the Jensen gap; a distribution narrow and centred near zero
implies a negligible gap.

\subsection{First-order approximation}

Since $F$ is piecewise linear, let $\xstar(\pbar)$ denote the optimal
dispatch at $\pbar$ (an element of the subdifferential $\partial F(\pbar)$).
The tangent approximation gives
\begin{equation}
  F(p_i)
  \;\approx\;
  F(\pbar) + \langle \xstar(\pbar),\, p_i - \pbar \rangle
  \quad
  \text{when $p_i$ and $\pbar$ lie in the same linear region of $F$.}
\end{equation}
Averaging over $i$:
\begin{equation}
  \E[F(p_i)]
  \;\approx\;
  F(\pbar) + \langle \xstar(\pbar),\, \underbrace{\E[p_i] - \pbar}_{=\,0} \rangle
  \;=\;
  F(\pbar).
\end{equation}
Hence $\Delta \approx 0$ whenever all scenarios remain in the same
linear region as $\pbar$.  The gap is non-trivial only if scenarios
\emph{cross a kink point} of $F$.

% ────────────────────────────────────────────────────────────────────────────
\section{Geometric Interpretation: Polytope Vertices}
% ────────────────────────────────────────────────────────────────────────────

\subsection{Finite vertex set}

The feasible set $\C$ is a polytope.  By the fundamental theorem of
linear programming, the optimum of any LP over $\C$ is attained at a
\emph{vertex}.  Therefore for each price scenario $p_i$ there exists
a vertex $v_i \in \mathrm{vert}(\C)$ such that
\begin{equation}
  F(p_i) \;=\; \langle c(p_i),\, v_i \rangle.
\end{equation}

\begin{proposition}[Finite vertex count]
The polytope $\C(s_0, s_1)$ has finitely many vertices.
For a 24-hour dispatch problem with $H = 24$ hours,
the number of vertices is bounded by $\binom{2H}{H}$ (the number of
basic feasible solutions of the underlying LP),
though in practice far fewer are active.
\end{proposition}

\subsection{Vertex dominance and the Jensen gap}

Let $\mathcal{V} = \{v^{(1)}, \ldots, v^{(K)}\}$ be the set of vertices
that are optimal for at least one scenario.  The value function can be
written as
\begin{equation}
  F(p) \;=\; \max_{k=1,\ldots,K} \langle c(p),\, v^{(k)} \rangle.
\end{equation}
Each $p \mapsto \langle c(p), v^{(k)} \rangle$ is linear.
The \emph{active vertex regions}
\begin{equation}
  R_k \;=\; \{p \in \R^{24} : v^{(k)} \text{ is optimal at } p\}
\end{equation}
partition (a.e.) the price space.

\begin{theorem}[Vertex dominance implies small gap]
  If all scenarios $p_1, \ldots, p_N$ lie in the same region $R_k$,
  then $F(p_i) = \langle c(p_i), v^{(k)} \rangle$ for all $i$, and
  \begin{equation}
    \Delta
    \;=\;
    \E[F(p_i)] - F(\pbar)
    \;=\;
    \bigl\langle c(\E[p_i]),\, v^{(k)} \bigr\rangle
    - \bigl\langle c(\pbar),\, v^{(k)} \bigr\rangle
    \;=\; 0.
  \end{equation}
\end{theorem}

\begin{remark}[Empirical finding]
  Numerical experiments with realistic German day-ahead price scenarios
  show that \textbf{the vast majority of scenarios activate the same vertex}
  as the mean price path.  The $k$-means clustering of dispatch profiles
  (Figure~\ref{fig:clusters}) reveals that one cluster accounts for
  over 90\% of scenarios on most days, confirming vertex dominance and
  explaining the near-zero extrinsic value observed in practice.
\end{remark}

% ────────────────────────────────────────────────────────────────────────────
\section{Cross-Sectional Analysis of \texorpdfstring{$F$}{F}}
% ────────────────────────────────────────────────────────────────────────────

\subsection{The normal direction}

At $\pbar$ the subdifferential of $F$ is $\partial F(\pbar) \ni \xstar(\pbar)$.
The \emph{zero-space} of the subdifferential is
\begin{equation}
  \ker(\xstar(\pbar))
  \;=\;
  \{d \in \R^{24} : \langle \xstar(\pbar), d \rangle = 0\},
\end{equation}
the set of price directions that do not change the optimal revenue to
first order.  The complementary direction — maximising the first-order
change — is
\begin{equation}
  v \;=\; \frac{\xstar(\pbar)}{\|\xstar(\pbar)\|}.
  \label{eq:vdir}
\end{equation}

\subsection{The 1-D cross-section}

We plot
\begin{equation}
  t \;\longmapsto\; F(\pbar + t\,v),
  \qquad
  t \in \bigl[-n_\sigma \cdot \sigma,\; n_\sigma \cdot \sigma\bigr],
\end{equation}
where $\sigma = \mathrm{std}(t_i)$ and $t_i = \langle p_i - \pbar, v\rangle$
are the scenario projections onto $v$.

The tangent at $t = 0$ is
\begin{equation}
  t \;\longmapsto\;
  F(\pbar) + t \cdot \langle \xstar(\pbar),\, v \rangle
  \;=\;
  F(\pbar) + t \cdot \|\xstar(\pbar)\|,
\end{equation}
since $\langle \xstar(\pbar), v \rangle = \|\xstar(\pbar)\|^2 / \|\xstar(\pbar)\| = \|\xstar(\pbar)\|$.

The \emph{convexity gap} $F(\pbar + tv) - \text{tangent}(t) \geq 0$
is the pointwise excess above the tangent.  It is zero within a linear
region and positive at kink crossings.

\begin{figure}[H]
  \centering
  \includegraphics[width=0.92\textwidth]{figures/cross_section.pdf}
  \caption{%
    Cross-section of $F$ along $v = \xstar(\pbar)/\|\xstar(\pbar)\|$
    for an example day.
    \textit{Top:} Blue curve — $F(\pbar + tv)$; navy dashed — tangent at $t=0$;
    red shading — convexity gap; orange dashed — vertex switch points.
    Secondary axis (tomato): gap $F - \text{tangent}$.
    \textit{Bottom:} Histogram of scenario projections $t_i$;
    orange dashed lines align with vertex switches in the upper panel.
    The scenarios cluster tightly around $t=0$, confirming small gap.
  }
  \label{fig:cross}
\end{figure}

\subsection{Key observation: asymmetric convexity}

The cross-section plots consistently show that $F$ is \textbf{strongly
convex for $t < 0$} (price moving against the optimal dispatch direction)
but \textbf{nearly linear for $t > 0$}.  This asymmetry arises because:

\begin{itemize}
  \item For $t > 0$: the current optimal dispatch $\xstar(\pbar)$ becomes
    even more profitable; the same vertex remains optimal and $F$ stays linear.
  \item For $t < 0$: the current dispatch becomes suboptimal; the LP must
    switch to a different vertex (different charge/discharge schedule),
    producing the observed kinks and curvature.
\end{itemize}

Since real price scenarios are centred near $\pbar$ and the convexity
appears predominantly at $t < 0$, the scenarios rarely reach the curved
region, and $\Delta \approx 0$.

% ────────────────────────────────────────────────────────────────────────────
\section{Dispatch Cluster Analysis}
% ────────────────────────────────────────────────────────────────────────────

To visualise vertex dominance directly, we apply $k$-means clustering
($k = 5$) to the 48-dimensional dispatch vectors
$(e^{\mathrm{in}}_i, e^{\mathrm{out}}_i) \in \R^{48}$
across all $N$ scenarios for a given day.

\begin{figure}[H]
  \centering
  \includegraphics[width=0.92\textwidth]{figures/clusters.pdf}
  \caption{%
    $k$-means cluster analysis of optimal dispatch profiles for an
    example day.  \textit{Top:} Cluster frequencies (\%).
    One cluster (C0) typically accounts for $> 90\%$ of scenarios.
    \textit{Remaining rows:} Mean dispatch shape (green = discharge,
    red = charge) and corresponding mean price shape per cluster.
    The dominant cluster corresponds to the same vertex as the
    intrinsic (mean-path) optimum.
  }
  \label{fig:clusters}
\end{figure}

The intra-cluster standard deviation $\sigma_k$ (annotated on each bar)
quantifies how close the LP solutions within a cluster are.
Small $\sigma_k$ confirms that each cluster genuinely corresponds to
a small neighbourhood of a single polytope vertex.

% ────────────────────────────────────────────────────────────────────────────
\section{Jensen Gap as a Function of Battery Parameters}
% ────────────────────────────────────────────────────────────────────────────

The shape of $F$ — and therefore the Jensen gap — depends on the
battery parameters through the feasible set $\C$.  Intuitively:

\begin{itemize}
  \item \textbf{Larger discharge power $P_d$}: more dispatch degrees of
    freedom; more vertices are reachable; kink points appear at smaller
    $|t|$; curvature increases; Jensen gap grows.
  \item \textbf{Lower charge efficiency $\eta_c$}: raises the effective
    charging cost; optimal dispatch shifts; the region around $\pbar$
    where the current vertex is optimal may shrink.
  \item \textbf{More energy levels}: finer SoC grid; richer continuation
    value; more accurate dynamic programme but does not directly affect
    the one-day convexity.
\end{itemize}

\begin{figure}[H]
  \centering
  \includegraphics[width=0.92\textwidth]{figures/param_sensitivity.pdf}
  \caption{%
    Convexity comparison across battery parameters.
    Two values of discharge power $P_d$ (left) and two values of
    charging efficiency $\eta_c$ (right), each shown at full scale and
    zoomed to the $\pm 2\sigma$ scenario range.
    Larger $P_d$ or lower $\eta_c$ produces a wider and earlier-appearing
    convexity gap, increasing the potential extrinsic value.
  }
  \label{fig:params}
\end{figure}

% ────────────────────────────────────────────────────────────────────────────
\section{Value Function over the Energy Grid}
% ────────────────────────────────────────────────────────────────────────────

The LSMC backward induction computes the value function
$V : \{s_{\min}, \ldots, s_{\max}\} \times \{1,\ldots,N\} \times \{0,\ldots,D{-}1\} \to \R$
on a discrete energy grid.  For a fixed day $d$:

\begin{equation}
  V(s, \omega, d)
  \;=\;
  \max_{s' \in \text{grid},\; x \in \C(s, s')}
  \Bigl[
    \langle c(p^\omega_d), x \rangle
    +
    \hat{V}(s', d+1; p^\omega_d)
  \Bigr],
\end{equation}
where $\hat{V}(\cdot, d+1; p)$ is the regressed continuation value
obtained by least-squares regression of $V(\cdot, \cdot, d+1)$ on a
polynomial basis in the mean daily price.

\begin{figure}[H]
  \centering
  \includegraphics[width=0.75\textwidth]{figures/value_dist.pdf}
  \caption{%
    Value function $V(s, \cdot, d)$ for an example day as a function of
    starting energy level $s$.
    Mean (blue curve) with P5--P95 (light) and P10--P90 (dark) scenario
    quantile bands.  Tomato dashed line: intrinsic value-to-go
    $V^{\mathrm{intr}}(d) = V^{\mathrm{intr}}_{\mathrm{total}} - \sum_{d'<d} r_{d'}^{\mathrm{intr}}$.
    The positive slope is the marginal EUR value of storing one extra MWh
    at this day boundary.
  }
  \label{fig:vdist}
\end{figure}

% ────────────────────────────────────────────────────────────────────────────
\section{Summary of Findings}
% ────────────────────────────────────────────────────────────────────────────

\begin{enumerate}
  \item \textbf{$F$ is convex and piecewise linear.}
    This is a mathematical consequence of the LP structure, not a modelling
    assumption.  The Jensen inequality $\E[F(p_i)] \geq F(\pbar)$ therefore
    holds with equality iff all scenarios lie in the same linear region.

  \item \textbf{The Jensen gap is the extrinsic value for one day.}
    The full (stochastic) value decomposes as
    $V^{\mathrm{full}} = V^{\mathrm{intr}} + \Delta$,
    where $\Delta \geq 0$ is the Jensen gap summed over all days.

  \item \textbf{Vertex dominance explains a near-zero extrinsic value.}
    Empirically, $> 90\%$ of scenarios activate the same LP vertex as the
    mean price path.  When all scenarios share a vertex, $F$ is linear
    on the scenario cloud and $\Delta = 0$ exactly.

  \item \textbf{The convexity is asymmetric: strong for $t < 0$, weak for $t > 0$.}
    Moving prices against the optimal dispatch (reducing spreads) forces
    a vertex switch and creates curvature.  Moving prices in the favourable
    direction keeps the current vertex and leaves $F$ linear.

  \item \textbf{Wider kink regions arise with larger discharge capacity or lower efficiency.}
    Batteries with more aggressive dispatch capabilities have more reachable
    vertices, producing kinks at smaller deviations from $\pbar$ and
    therefore larger potential Jensen gaps.

  \item \textbf{The slope of $V(s, \cdot, d)$ measures marginal storage value.}
    A steep positive slope indicates that an extra MWh of stored energy
    at day boundary $d$ is worth carrying into the next day.
    A flat $V$ implies the end-of-day SoC level does not matter for
    continuation value.
\end{enumerate}

% ────────────────────────────────────────────────────────────────────────────
\appendix
\section{LP Formulation}
\label{app:lp}
% ────────────────────────────────────────────────────────────────────────────

For horizon $H = 24$ hours, start energy $s_0$, end energy $s_1$,
charge power limit $P_c$, discharge power limit $P_d$,
time step $\Delta t$, min/max energy $[s_{\min}, s_{\max}]$:

\textbf{Variables:} $e^{\mathrm{in}} \in \R^H$, $e^{\mathrm{out}} \in \R^H$.

\textbf{Objective (maximise):}
\begin{equation}
  \sum_{h=0}^{H-1}
  \Bigl(
    p_h \cdot \eta_d \cdot e^{\mathrm{out}}_h
    \;-\;
    \frac{p_h}{\eta_c} \cdot e^{\mathrm{in}}_h
  \Bigr).
\end{equation}

\textbf{Box constraints:}
$0 \leq e^{\mathrm{in}}_h \leq P_c \Delta t$,\quad
$0 \leq e^{\mathrm{out}}_h \leq P_d \Delta t$.

\textbf{Energy balance:}
$\displaystyle\sum_{h} e^{\mathrm{in}}_h - \sum_{h} e^{\mathrm{out}}_h = s_1 - s_0$.

\textbf{Intra-day path constraints} (triangular matrix $L = \mathrm{tril}(\mathbf{1})$):
\begin{align}
  (Le^{\mathrm{in}} - Le^{\mathrm{out}})_h &\leq s_{\max} - s_0, \\
  (Le^{\mathrm{out}} - Le^{\mathrm{in}})_h &\leq s_0 - s_{\min},
  \qquad h = 0, \ldots, H{-}1.
\end{align}

These constraints ensure the running energy level
$s_h = s_0 + \sum_{h'=0}^{h-1}(e^{\mathrm{in}}_{h'} - e^{\mathrm{out}}_{h'})$
stays within $[s_{\min}, s_{\max}]$ at every hour.

\end{document}
